% Puzzle difficulty without the solution: the paper. The prose follows
% docs/writing.md. Every number, table, figure and board is written by
% `lab paper puzzle-difficulty-unsolved` into generated/, from the model's
% committed results and from those of the model it is compared with, and read
% back with \V{key} and \block{name}.
\documentclass[11pt]{article}
% The preamble every paper of the lab shares: % The preamble every paper of the lab shares: % The preamble every paper of the lab shares: \input{../../../papers/preamble}
% from models/<model>/paper/, \input{../preamble} from papers/<name>/.
% Boards come from chessboard and xskak, plots from pgfplots. A paper prints no
% number by hand: `lab paper` writes every number, table and figure under the
% paper's generated/ folder, and paper.tex says what they mean.
\usepackage[LSBC4,T1]{fontenc}
\usepackage[utf8]{inputenc}
\usepackage{lmodern}
\usepackage[margin=1in]{geometry}
\usepackage{microtype}
\usepackage{amsmath}
\usepackage{booktabs}
\usepackage{tabularx}
\usepackage{array}
\usepackage{xcolor}
\usepackage[obeyspaces,spaces]{url}
\usepackage{xskak}
\usepackage{pgfplots}
\usepgfplotslibrary{groupplots}
\usetikzlibrary{positioning,arrows.meta,calc,fit}
\pgfplotsset{compat=1.18}
\usepackage{caption}
\usepackage{subcaption}
\usepackage{enumitem}
\usepackage[section]{placeins}
\usepackage[hidelinks]{hyperref}

% A generated value by its key: generated/values.tex sets them, \V{key} prints
% one, and a key nobody set stops the build.
\makeatletter
\newcommand{\setvalue}[2]{\expandafter\gdef\csname labvalue@#1\endcsname{#2}}
\newcommand{\V}[1]{\@ifundefined{labvalue@#1}%
  {\PackageError{lab}{No generated value named '#1'}{Run `lab paper` for this model again.}}%
  {\csname labvalue@#1\endcsname}}
\makeatother

% Chart colours. A run keeps its colour in every figure of every paper.
\definecolor{labBlue}{HTML}{2A78D6}
\definecolor{labOrange}{HTML}{EB6834}
\definecolor{labAqua}{HTML}{1BAF7A}
\definecolor{labYellow}{HTML}{EDA100}
\definecolor{labMagenta}{HTML}{E87BA4}
\definecolor{labGreen}{HTML}{008300}
\definecolor{labViolet}{HTML}{6250D6}
\definecolor{labGray}{HTML}{898781}
\definecolor{labInk}{HTML}{0B0B0B}
\definecolor{labSecondary}{HTML}{52514E}
\definecolor{labGrid}{HTML}{E1E0D9}
\definecolor{labAxis}{HTML}{C3C2B7}
\definecolor{labWash}{HTML}{CDE2FB}

\pgfplotsset{
  lab/.style={
    font=\sffamily\footnotesize,
    axis line style={labAxis},
    tick style={labAxis},
    tick label style={labSecondary},
    label style={labSecondary},
    title style={labInk, font=\sffamily\small\bfseries, align=left},
    grid style={labGrid, line width=0.3pt},
    major grid style={labGrid, line width=0.3pt},
    axis lines*=left,
    tick align=outside,
    clip=false,
    legend style={draw=none, fill=none, font=\sffamily\footnotesize, text=labSecondary,
      cells={anchor=west}, legend columns=-1, column sep=6pt, at={(0,1.03)}, anchor=south west},
    every axis plot/.append style={line width=1.1pt, line join=round, line cap=round},
    scaled ticks=false,
    /pgf/number format/1000 sep={,},
  },
  lab marks/.style={mark=*, mark size=1.9pt, mark options={draw=white, line width=0.6pt}},
}

\setchessboard{boardfontencoding=LSBC4, boardfontfamily=skaknew, blackfieldcolor=black!18,
  whitefieldcolor=white, setfontcolors, showmover=false, label=true, labelfontsize=6pt,
  boardfontsize=12pt}

% A move or line in UCI, a position in FEN (breakable at a slash), and a run's name.
\newcommand{\uci}[1]{\texttt{#1}}
\newcommand{\fen}[1]{{\footnotesize\path{#1}}}
\newcommand{\run}[1]{\texttt{#1}}
\newcolumntype{L}{>{\raggedright\arraybackslash}X}
\newcolumntype{R}{>{\raggedleft\arraybackslash}X}

\setlength{\parskip}{0.4em}
\setlength{\parindent}{0pt}
\captionsetup{font=small, labelfont=bf}

% from models/<model>/paper/, \input{../preamble} from papers/<name>/.
% Boards come from chessboard and xskak, plots from pgfplots. A paper prints no
% number by hand: `lab paper` writes every number, table and figure under the
% paper's generated/ folder, and paper.tex says what they mean.
\usepackage[LSBC4,T1]{fontenc}
\usepackage[utf8]{inputenc}
\usepackage{lmodern}
\usepackage[margin=1in]{geometry}
\usepackage{microtype}
\usepackage{amsmath}
\usepackage{booktabs}
\usepackage{tabularx}
\usepackage{array}
\usepackage{xcolor}
\usepackage[obeyspaces,spaces]{url}
\usepackage{xskak}
\usepackage{pgfplots}
\usepgfplotslibrary{groupplots}
\usetikzlibrary{positioning,arrows.meta,calc,fit}
\pgfplotsset{compat=1.18}
\usepackage{caption}
\usepackage{subcaption}
\usepackage{enumitem}
\usepackage[section]{placeins}
\usepackage[hidelinks]{hyperref}

% A generated value by its key: generated/values.tex sets them, \V{key} prints
% one, and a key nobody set stops the build.
\makeatletter
\newcommand{\setvalue}[2]{\expandafter\gdef\csname labvalue@#1\endcsname{#2}}
\newcommand{\V}[1]{\@ifundefined{labvalue@#1}%
  {\PackageError{lab}{No generated value named '#1'}{Run `lab paper` for this model again.}}%
  {\csname labvalue@#1\endcsname}}
\makeatother

% Chart colours. A run keeps its colour in every figure of every paper.
\definecolor{labBlue}{HTML}{2A78D6}
\definecolor{labOrange}{HTML}{EB6834}
\definecolor{labAqua}{HTML}{1BAF7A}
\definecolor{labYellow}{HTML}{EDA100}
\definecolor{labMagenta}{HTML}{E87BA4}
\definecolor{labGreen}{HTML}{008300}
\definecolor{labViolet}{HTML}{6250D6}
\definecolor{labGray}{HTML}{898781}
\definecolor{labInk}{HTML}{0B0B0B}
\definecolor{labSecondary}{HTML}{52514E}
\definecolor{labGrid}{HTML}{E1E0D9}
\definecolor{labAxis}{HTML}{C3C2B7}
\definecolor{labWash}{HTML}{CDE2FB}

\pgfplotsset{
  lab/.style={
    font=\sffamily\footnotesize,
    axis line style={labAxis},
    tick style={labAxis},
    tick label style={labSecondary},
    label style={labSecondary},
    title style={labInk, font=\sffamily\small\bfseries, align=left},
    grid style={labGrid, line width=0.3pt},
    major grid style={labGrid, line width=0.3pt},
    axis lines*=left,
    tick align=outside,
    clip=false,
    legend style={draw=none, fill=none, font=\sffamily\footnotesize, text=labSecondary,
      cells={anchor=west}, legend columns=-1, column sep=6pt, at={(0,1.03)}, anchor=south west},
    every axis plot/.append style={line width=1.1pt, line join=round, line cap=round},
    scaled ticks=false,
    /pgf/number format/1000 sep={,},
  },
  lab marks/.style={mark=*, mark size=1.9pt, mark options={draw=white, line width=0.6pt}},
}

\setchessboard{boardfontencoding=LSBC4, boardfontfamily=skaknew, blackfieldcolor=black!18,
  whitefieldcolor=white, setfontcolors, showmover=false, label=true, labelfontsize=6pt,
  boardfontsize=12pt}

% A move or line in UCI, a position in FEN (breakable at a slash), and a run's name.
\newcommand{\uci}[1]{\texttt{#1}}
\newcommand{\fen}[1]{{\footnotesize\path{#1}}}
\newcommand{\run}[1]{\texttt{#1}}
\newcolumntype{L}{>{\raggedright\arraybackslash}X}
\newcolumntype{R}{>{\raggedleft\arraybackslash}X}

\setlength{\parskip}{0.4em}
\setlength{\parindent}{0pt}
\captionsetup{font=small, labelfont=bf}

% from models/<model>/paper/, \input{../preamble} from papers/<name>/.
% Boards come from chessboard and xskak, plots from pgfplots. A paper prints no
% number by hand: `lab paper` writes every number, table and figure under the
% paper's generated/ folder, and paper.tex says what they mean.
\usepackage[LSBC4,T1]{fontenc}
\usepackage[utf8]{inputenc}
\usepackage{lmodern}
\usepackage[margin=1in]{geometry}
\usepackage{microtype}
\usepackage{amsmath}
\usepackage{booktabs}
\usepackage{tabularx}
\usepackage{array}
\usepackage{xcolor}
\usepackage[obeyspaces,spaces]{url}
\usepackage{xskak}
\usepackage{pgfplots}
\usepgfplotslibrary{groupplots}
\usetikzlibrary{positioning,arrows.meta,calc,fit}
\pgfplotsset{compat=1.18}
\usepackage{caption}
\usepackage{subcaption}
\usepackage{enumitem}
\usepackage[section]{placeins}
\usepackage[hidelinks]{hyperref}

% A generated value by its key: generated/values.tex sets them, \V{key} prints
% one, and a key nobody set stops the build.
\makeatletter
\newcommand{\setvalue}[2]{\expandafter\gdef\csname labvalue@#1\endcsname{#2}}
\newcommand{\V}[1]{\@ifundefined{labvalue@#1}%
  {\PackageError{lab}{No generated value named '#1'}{Run `lab paper` for this model again.}}%
  {\csname labvalue@#1\endcsname}}
\makeatother

% Chart colours. A run keeps its colour in every figure of every paper.
\definecolor{labBlue}{HTML}{2A78D6}
\definecolor{labOrange}{HTML}{EB6834}
\definecolor{labAqua}{HTML}{1BAF7A}
\definecolor{labYellow}{HTML}{EDA100}
\definecolor{labMagenta}{HTML}{E87BA4}
\definecolor{labGreen}{HTML}{008300}
\definecolor{labViolet}{HTML}{6250D6}
\definecolor{labGray}{HTML}{898781}
\definecolor{labInk}{HTML}{0B0B0B}
\definecolor{labSecondary}{HTML}{52514E}
\definecolor{labGrid}{HTML}{E1E0D9}
\definecolor{labAxis}{HTML}{C3C2B7}
\definecolor{labWash}{HTML}{CDE2FB}

\pgfplotsset{
  lab/.style={
    font=\sffamily\footnotesize,
    axis line style={labAxis},
    tick style={labAxis},
    tick label style={labSecondary},
    label style={labSecondary},
    title style={labInk, font=\sffamily\small\bfseries, align=left},
    grid style={labGrid, line width=0.3pt},
    major grid style={labGrid, line width=0.3pt},
    axis lines*=left,
    tick align=outside,
    clip=false,
    legend style={draw=none, fill=none, font=\sffamily\footnotesize, text=labSecondary,
      cells={anchor=west}, legend columns=-1, column sep=6pt, at={(0,1.03)}, anchor=south west},
    every axis plot/.append style={line width=1.1pt, line join=round, line cap=round},
    scaled ticks=false,
    /pgf/number format/1000 sep={,},
  },
  lab marks/.style={mark=*, mark size=1.9pt, mark options={draw=white, line width=0.6pt}},
}

\setchessboard{boardfontencoding=LSBC4, boardfontfamily=skaknew, blackfieldcolor=black!18,
  whitefieldcolor=white, setfontcolors, showmover=false, label=true, labelfontsize=6pt,
  boardfontsize=12pt}

% A move or line in UCI, a position in FEN (breakable at a slash), and a run's name.
\newcommand{\uci}[1]{\texttt{#1}}
\newcommand{\fen}[1]{{\footnotesize\path{#1}}}
\newcommand{\run}[1]{\texttt{#1}}
\newcolumntype{L}{>{\raggedright\arraybackslash}X}
\newcolumntype{R}{>{\raggedleft\arraybackslash}X}

\setlength{\parskip}{0.4em}
\setlength{\parindent}{0pt}
\captionsetup{font=small, labelfont=bf}

% Written by `lab paper puzzle-difficulty-unsolved`; do not edit by hand.
\looktrue
\chancefalse
\pairstrue
\variantsfalse
\gamesfalse
\featuresfalse
\changestrue
\showntrue
\shownpairedfalse
\setvalue{bands.best-ahead.of}{three of the five bands}
\setvalue{bands.count}{five}
\setvalue{bands.lead.narrowest}{$-$29.5}
\setvalue{bands.lead.narrowest.name}{1000-1399}
\setvalue{bands.lead.widest}{+122}
\setvalue{bands.lead.widest.name}{\textless{}1000}
\setvalue{best.close.median}{267.3}
\setvalue{best.close.p90}{678.6}
\setvalue{best.close.within.100}{20}
\setvalue{best.close.within.200}{39}
\setvalue{best.close.within.300}{55}
\setvalue{best.mae}{324.3}
\setvalue{best.name}{\run{conv-small}}
\setvalue{best.parameters.millions}{1.0}
\setvalue{best.r2}{0.3294}
\setvalue{best.reliability.groups}{20}
\setvalue{best.reliability.highest.happened}{1938}
\setvalue{best.reliability.highest.said}{1991}
\setvalue{best.reliability.lowest.happened}{737}
\setvalue{best.reliability.lowest.said}{708}
\setvalue{best.reliability.spread}{1120}
\setvalue{best.reliability.widest-gap}{52.2}
\setvalue{best.score}{411.8}
\setvalue{best.score.high}{413.6}
\setvalue{best.score.low}{409.9}
\setvalue{best.seconds}{3 min}
\setvalue{best.spearman}{0.5605}
\setvalue{challenger.beside.closer}{57}
\setvalue{challenger.beside.correlation}{0.57}
\setvalue{challenger.name}{\run{conv-small}}
\setvalue{code.commit}{\texttt{0ba155a}}
\setvalue{cost.seconds}{34 min}
\setvalue{cost.seconds.networks}{33 min}
\setvalue{cost.seconds.others}{34 s}
\setvalue{data.fingerprint}{\texttt{d14731fa4425}}
\setvalue{data.licence}{CC0 1.0}
\setvalue{data.name}{\texttt{puzzles-v1}}
\setvalue{data.rows}{1{,}127{,}915}
\setvalue{data.rows.test}{112{,}450}
\setvalue{data.share.test}{10}
\setvalue{data.share.train}{85}
\setvalue{data.share.valid}{5}
\setvalue{examples.count}{three}
\setvalue{floor.constant}{1462.4}
\setvalue{floor.score}{502.9}
\setvalue{grid.smallest-cell}{30}
\setvalue{label.mean}{1462.4}
\setvalue{label.q5}{693.0}
\setvalue{label.q95}{2338.0}
\setvalue{label.sd}{503.3}
\setvalue{ladder.at-last-epoch.of}{three of the four networks}
\setvalue{look.date}{22 November 2026}
\setvalue{look.resamples}{200}
\setvalue{look.smallest-compared}{200}
\setvalue{look.test-rows}{112{,}450}
\setvalue{method.gbm.features.percent}{80}
\setvalue{method.gbm.leaves}{127}
\setvalue{method.gbm.patience}{100}
\setvalue{method.gbm.rate}{0.05}
\setvalue{method.gbm.rounds}{4{,}000}
\setvalue{method.gbm.rows.percent}{80}
\setvalue{method.resamples}{1{,}000}
\setvalue{networks.ahead.of}{all four networks}
\setvalue{networks.patience}{three}
\setvalue{networks.warmup.percent}{5}
\setvalue{pair.transformer-medium-hand.transformer-medium.gain}{+9.61}
\setvalue{pair.transformer-medium-hand.transformer-medium.valid.verdict}{better}
\setvalue{pair.transformer-medium-hand.transformer-medium.verdict}{better beyond the noise}
\setvalue{recipe.filters.max-rating-deviation}{100}
\setvalue{recipe.filters.min-plays}{100}
\setvalue{recipe.sample.percent}{25}
\setvalue{reference.features}{18}
\setvalue{reference.leans.1.feature}{checks on offer}
\setvalue{reference.leans.1.share}{19}
\setvalue{reference.leans.2.feature}{material on the board}
\setvalue{reference.leans.2.share}{13}
\setvalue{reference.leans.3.feature}{material balance}
\setvalue{reference.leans.3.share}{12}
\setvalue{reference.lift.percent}{9.2}
\setvalue{reference.score}{456.8}
\setvalue{reference.trees}{1{,}212}
\setvalue{results.date}{22 November 2026}
\setvalue{run.conv-small.gain}{+45.01}
\setvalue{run.transformer-medium-hand.gain}{+40.71}
\setvalue{run.transformer-medium-hand.name}{\run{transformer-medium-hand}}
\setvalue{run.transformer-medium-hand.parameters.millions}{6.6}
\setvalue{run.transformer-medium.gain}{+31.10}
\setvalue{run.transformer-medium.name}{\run{transformer-medium}}
\setvalue{run.transformer-small.gain}{+27.11}
\setvalue{runs.count}{six}
\setvalue{shown.best.here}{411.8}
\setvalue{shown.best.here.name}{\run{conv-small}}
\setvalue{shown.best.kept.percent}{38}
\setvalue{shown.best.lost.percent}{62}
\setvalue{shown.best.there}{261.9}
\setvalue{shown.best.there.name}{\run{transformer-medium-hand-long}}
\setvalue{shown.best.worth}{150}
\setvalue{shown.model}{puzzle difficulty}
\setvalue{shown.networks.least.name}{\run{conv-small}}
\setvalue{shown.networks.least.worth}{121}
\setvalue{shown.networks.most.name}{\run{transformer-medium-hand}}
\setvalue{shown.networks.most.worth}{146}
\setvalue{shown.reference.here}{456.8}
\setvalue{shown.reference.here.lift.percent}{9.2}
\setvalue{shown.reference.kept.percent}{22}
\setvalue{shown.reference.there}{294.9}
\setvalue{shown.reference.there.lift.percent}{41}
\setvalue{shown.reference.worth}{162}
\setvalue{shown.runs}{six}
\setvalue{shown.there.conv-small.gain}{+4.105}
\setvalue{shown.there.gbm-context.lift.percent}{15}
\setvalue{shown.there.gbm-context.name}{\run{gbm-context}}
\setvalue{shown.there.gbm-context.over-best}{16.9}
\setvalue{shown.there.gbm-context.over-reference}{28.1}
\setvalue{shown.there.gbm-context.score}{428.8}
\setvalue{shown.there.gbm-context.than-best}{behind}
\setvalue{shown.there.gbm-context.than-reference}{ahead of}
\setvalue{shown.there.transformer-medium-hand.gain}{+24.68}
\setvalue{shown.there.transformer-medium.gain}{+10.91}
\setvalue{shown.there.transformer-small.gain}{+1.847}
\setvalue{shown.transformer-medium-hand.kept.percent}{37}
\setvalue{shown.transformer-medium.kept.percent}{35}
\setvalue{slice.band.-1000.challenger.off}{+327}
\setvalue{slice.band.-1000.reference.off}{+508}
\setvalue{slice.band.2200+.challenger.off}{$-$720}
\setvalue{slice.band.2200+.reference.off}{$-$826}
\setvalue{slice.check.no.challenger.off}{$-$99}
\setvalue{slice.check.no.happened}{1724}
\setvalue{slice.check.yes.challenger.off}{+61}
\setvalue{slice.check.yes.happened}{1300}
\setvalue{slice.mate.no.challenger.off}{$-$53}
\setvalue{slice.mate.no.happened}{1613}
\setvalue{slice.mate.yes.challenger.off}{+120}
\setvalue{slice.mate.yes.happened}{1119}
\setvalue{slice.moves.1.challenger.off}{+189}
\setvalue{slice.moves.1.happened}{961}
\setvalue{slice.moves.1.reference.off}{+393}
\setvalue{slice.moves.2.happened}{1395}
\setvalue{slice.moves.3.happened}{1688}
\setvalue{slice.moves.5+.challenger.off}{$-$469}
\setvalue{slice.moves.5+.happened}{2098}
\setvalue{slice.moves.5+.reference.off}{$-$534}
\setvalue{slice.quiet.0.challenger.off}{+86}
\setvalue{slice.quiet.0.happened}{1301}
\setvalue{slice.quiet.1.challenger.off}{$-$155}
\setvalue{slice.quiet.1.happened}{1759}
\setvalue{slice.sacrifice.no.challenger.off}{+6.4}
\setvalue{slice.sacrifice.no.happened}{1443}
\setvalue{slice.sacrifice.yes.challenger.off}{$-$40}
\setvalue{slice.sacrifice.yes.happened}{1580}
\setvalue{slice.themes.ahead}{30}
\setvalue{slice.themes.behind}{zero}
\setvalue{slice.themes.compared}{30}
\setvalue{slice.themes.gain.least}{+13.79}
\setvalue{slice.themes.gain.least.group}{promotion}
\setvalue{slice.themes.gain.most}{+186.4}
\setvalue{slice.themes.gain.most.group}{backRankMate}
\setvalue{task.context.count}{four}
\setvalue{task.hand.count}{14}
\setvalue{task.marks.count}{three}
\setvalue{task.reference}{\run{gbm-hand}}


\begin{document}

\papertitle{Predicting the Difficulty of Chess Puzzles from the Position
  Alone}
  {Results as of \V{results.date}\quad\textperiodcentered\quad dataset \V{data.name},
   build \V{data.fingerprint}\quad\textperiodcentered\quad code \V{code.commit}}

\begin{abstract}
A model that is given the solution of a chess puzzle can predict the puzzle's
Lichess rating from properties of the solution, such as its length. We ask how
much of the rating can be predicted from the position alone: the board
presented to the solver and the opponent's preceding move, with all
information about the solution withheld. The dataset, the split and the model
configurations are those of a companion study in which the models are given
the solution, so each model has a counterpart that differs from it only in
its inputs. On \V{data.rows.test} test puzzles, the best of \V{runs.count}
models is a transformer that also receives hand-crafted features of the
position and whose encoder is pretrained on engine evaluations. Its
root-mean-square error (RMSE) is \V{best.score} rating points (95\% CI
\V{best.score.low} to \V{best.score.high}), compared with
\V{reference.score} for gradient-boosted decision trees (GBDT) on the
hand-crafted features and \V{floor.score} for a constant baseline.
\ifshown
Its counterpart with the solution has an RMSE of \V{shown.best.there}, so
without the solution the model achieves \V{shown.best.kept.percent}\% of the
counterpart's improvement over the constant baseline. GBDT achieves
\V{shown.reference.kept.percent}\%.
\fi
Pretraining is the most effective modification that we tested. The
predictions order puzzles moderately well (Spearman correlation
\V{best.spearman}) and regress strongly towards the mean rating, so they are
not suitable for rating an individual position.
\end{abstract}

\section{Introduction}
\label{sec:intro}

A Lichess puzzle is solved from the position. The solver sees the board and
the move the opponent has just played, knows that a winning line exists, and
has to find it~\cite{lichessdb}. The puzzle's Glicko-2
rating~\cite{glicko2} summarises the outcomes of many such attempts, and it
is the best available measure of how difficult the tactic is for human
solvers.

A model can predict this rating in two settings. In the first, the model is
given the solution, in the way an analyst who has an engine's line asks how
hard that line was to find. The companion paper, \emph{Predicting the
Difficulty of Chess Puzzles from the Position and Its Solution}, studies this
setting and finds that hand-crafted features of the solution carry most of
the signal. In the second setting, the model receives only what the solver
sees. This paper studies the second setting on the same puzzles.

We address two questions. The first is how accurately the difficulty of a
puzzle can be predicted from its position. The second is how much of the
accuracy of a model that is given the solution depends on the solution. The
second question is answered by a controlled comparison: every model in this
paper has a counterpart in the companion study with the same configuration,
the same training puzzles and the same test puzzles, and the two differ only
in whether the solution is part of the input.

This work supports a chess coaching application, for which the comparison
shows where a difficulty model has to be placed. If most of a rating could be
predicted from the position, positions could be ordered by difficulty before
any line is computed. If most of it depends on the solution, the difficulty
model has to follow the engine analysis that provides the line. Our main
findings are the following.
\begin{itemize}
\item \textbf{Aggregate features of the position predict little.} The GBDT
baseline, which receives \V{task.hand.count} hand-crafted features of the
position, has an RMSE of \V{reference.score} rating points, only
\V{reference.lift.percent}\% below the constant baseline (\V{floor.score}).
\item \textbf{Neural networks that read the board predict much more.}
Compared with the GBDT baseline, \V{networks.ahead.of} are significantly
better. The best model has a gain of \V{best.gain} RMSE points (95\% CI
\V{best.gain.low} to \V{best.gain.high}), which is \V{best.gain.of-lift}\% of
the improvement of the GBDT baseline over the constant baseline
(Section~\ref{sec:results}).
\ifshown
\item \textbf{About half of the improvement is achieved without the
solution.} The best model achieves \V{shown.best.kept.percent}\% of the
improvement over the constant baseline that its counterpart achieves with
the solution. The medium transformer trained from random initialisation
achieves \V{shown.transformer-medium.kept.percent}\% and GBDT
\V{shown.reference.kept.percent}\% (Section~\ref{sec:shown}).
\fi
\ifchangestart
\item \textbf{Pretraining on engine evaluations is the most effective
modification.} It improves the medium transformer by
\V{pair.transformer-medium-pretrained.transformer-medium.gain} RMSE points
(95\% CI \V{pair.transformer-medium-pretrained.transformer-medium.gain.low}
to \V{pair.transformer-medium-pretrained.transformer-medium.gain.high}),
several times the gain from adding the hand-crafted features
(Section~\ref{sec:ablations}).
\fi
\iflook
\item \textbf{The predictions regress strongly towards the mean.} For puzzles
rated 2200 and above, the mean prediction of the best model minus the mean
true rating is \V{slice.band.2200+.challenger.off} points, and for puzzles
rated below 1000 it is \V{slice.band.-1000.challenger.off} points. The
models recover only part of the solution length, a property that by itself
predicts the rating better than all hand-crafted features of the position
(Sections~\ref{sec:shown} and~\ref{sec:error}).
\fi
\end{itemize}

\section{Related Work}
\label{sec:related}

\paragraph{Puzzle difficulty.}
Predicting the Lichess rating of a puzzle from its position and solution was
the task of data-mining competitions at IEEE BigData 2024 and FedCSIS
2025~\cite{puzzlecup}. The winning entry of the 2025 competition combines a
fine-tuned Maia-2~\cite{maia2}, a model of human moves, with hand-crafted and
engine-derived features~\cite{puzzlecupwinner}. Mi{\l}osz and
Kapusta~\cite{glickformer} introduce GlickFormer, a transformer that encodes
the sequence of positions along the solution and predicts the puzzle's
Glicko-2 rating. In these approaches the model is given the solution. We
withhold it, in order to measure how much of the rating can be predicted from
the position that the solver sees, and we compare each model with a
counterpart that is given the solution. We use our own sample of the puzzle
database and our own test set, so our scores are not comparable with the
competition results.

\paragraph{Models of human play.}
Maia~\cite{maia} and Maia-2~\cite{maia2} predict the moves that players of a
given rating make. Anderson et al.~\cite{humanerror} find that the difficulty
of a position predicts human error better than the player's skill or the time
available. These studies model behaviour in games. Puzzle ratings offer a
complementary signal: a direct measurement, from many attempts, of how hard
one specific tactic is.

\paragraph{Learning from engine evaluations.}
Ruoss et al.~\cite{searchless} train transformers to predict engine
evaluations and play strong chess without search. Our pretrained encoder is
trained on the same kind of target at a much smaller scale, and we use it
only as an initialisation for the difficulty model.

\section{Task and Data}
\label{sec:data}

\subsection{Task}

The input is the position that the solver sees and the opponent's preceding
move. The target is the puzzle's Lichess rating. Table~\ref{tab:brief}
summarises the task.

\begin{table}[htbp]\centering\small
\begin{tabularx}{\textwidth}{lL}
\toprule
Input & The position presented to the solver and the opponent's preceding
move, whose origin and destination squares are marked in auxiliary input
planes. No information about the solution. \\
Target & The puzzle's Lichess rating, in rating points. \\
Data & The Lichess puzzle database (\V{data.licence}): a fixed
\V{recipe.sample.percent}\% sample of the puzzles with a stable rating, the
same dataset as in the companion study. \\
Split & By the game a puzzle was extracted from: \V{data.share.train}\%
training, \V{data.share.valid}\% validation, \V{data.share.test}\% test, with
the same puzzles in each set as in the companion study. \\
Primary metric & RMSE on the test set, in rating points (lower is better). \\
Comparison & The model of the same name in the companion study, which differs
only in that it is given the solution. \\
\bottomrule
\end{tabularx}
\caption{Summary of the task.}
\label{tab:brief}
\end{table}

\subsection{Dataset}

The source is the Lichess puzzle database~\cite{lichessdb}. Each record gives
the puzzle's identifier, the position before the opponent's move, the moves
of the solution with the opponent's move first, the puzzle's rating and
rating deviation, the number of attempts, the theme tags, and the game the
puzzle was extracted from. The solution is thus part of every record. In this
task it is never an input, and we use it only to analyse the errors of the
models (Section~\ref{sec:error}).

We keep the puzzles whose rating is stable: a Glicko-2 rating deviation of at
most \V{recipe.filters.max-rating-deviation} after at least
\V{recipe.filters.min-plays} attempts. Of these we keep a fixed
\V{recipe.sample.percent}\% sample, selected by a hash of the puzzle
identifier, so that the dataset can be rebuilt exactly. The resulting dataset
has \V{data.rows} puzzles (Table~\ref{tab:data}).

The split is determined by a hash of the source game, so no game contributes
puzzles to more than one of the training, validation and test sets. The same
grouping is used for the bootstrap (Section~\ref{sec:protocol}). The dataset
is not stored in the repository. Its recipe is, and a build of the dataset is
identified by a fingerprint, a hash of its contents. All models in this
paper and their counterparts in the companion study were trained on build
\V{data.fingerprint}, and the evaluation code refuses to compare models
trained on different builds.

\begin{table}[htbp]\centering\small
\block{data}
\caption{Dataset statistics by split.}
\label{tab:data}
\end{table}

\iflook
The source contains \V{manifest.source-rows-read} puzzles, of which
\V{manifest.kept.percent}\% remain after filtering and sampling. The ratings
have a mean of \V{label.mean} and a standard deviation of \V{label.sd}, and
90\% of them lie between \V{label.q5} and \V{label.q95}
(Figure~\ref{fig:label}, Table~\ref{tab:describe}). The standard deviation is
also the RMSE of a model that predicts the mean rating for every puzzle.

\begin{figure}[htbp]\centering
\block{label}
\caption{Distribution of the puzzle ratings in the dataset.}
\label{fig:label}
\end{figure}

\begin{table}[htbp]\centering\small
\block{describe}
\caption{Summary statistics of selected variables over the whole dataset:
the mean, five quantiles, the minimum and the maximum. The solution length is
listed for reference and is not an input.}
\label{tab:describe}
\end{table}
\fi

\section{Methods}
\label{sec:methods}

\subsection{Input representation}
\label{sec:input}

% Input representation: identical for every task, so each model paper inputs
% this file under its own heading and then adds what is specific to its task.
% The prose follows docs/writing.md.
All tasks in this project share one input representation, so two tasks differ
in their target and in the information a model is given, and not in how a
position is encoded.

\paragraph{Square tokens.}
A position is encoded as 64 tokens, one per square. Token 0 denotes an empty
square, tokens 1 to 6 the pawn, knight, bishop, rook, queen and king of the
side to move, which we call the \emph{mover}, and tokens 7 to 12 the
opponent's pieces. When Black is to move, the board is flipped vertically and
the colours are swapped, so that the mover always plays up the board. A
position with White to move and the corresponding colour-reversed position
with Black to move therefore have the same encoding, and a model does not
have to learn the two cases separately.

\paragraph{Auxiliary input planes.}
Information about individual squares is supplied as binary planes over the 64
squares, which a network receives together with the tokens. This task uses
\V{task.marks.count} planes (Table~\ref{tab:marks}). The first marks the en
passant target square, and the others mark the origin and destination squares
of moves. A plane identifies the squares of a move and carries no other
information about it, such as whether the move is a capture or a check.

\begin{table}[htbp]\centering\small
\block{marks}
\caption{Auxiliary input planes. Squares are given from the mover's point of
view, after the board has been flipped.}
\label{tab:marks}
\end{table}

\ifsequence
\paragraph{Sequences of positions.}
An example of this task is not a single position. It is the sequence of one
player's moves in a game, and each element of the sequence is a position
encoded as above: the board before the move, with the opponent's preceding
move and the move played marked in the auxiliary planes. Scalar inputs
therefore come at two levels, one for each move and one for the game.

\paragraph{Per-move features.}
Each position of the sequence has \V{task.steps.count} scalar features of its
own (Table~\ref{tab:steps}): the four castling rights and the task's
information about that move. A network receives them together with the
position they belong to.

\begin{table}[htbp]\centering\small
\block{steps}
\caption{Per-move features: the scalar inputs attached to each position of a
sequence.}
\label{tab:steps}
\end{table}

\paragraph{Scalar context features.}
Information about the game as a whole is given as \V{task.context.count}
scalar context features (Table~\ref{tab:reads}). Every model receives them,
including the gradient-boosted trees, which never see a board. For a network,
both kinds of scalar feature are standardised with the mean and standard
deviation of the training set.

\begin{table}[htbp]\centering\small
\block{reads}
\caption{Scalar context features: the game-level inputs given to every
model.}
\label{tab:reads}
\end{table}

\paragraph{Hand-crafted features.}
Gradient-boosted trees cannot take a sequence of boards as input, so we also
compute \V{task.hand.count} aggregate features over the player's moves
(Table~\ref{tab:hand} in the appendix): how much win probability the moves
lost on average, at worst and how often beyond fixed thresholds, how the clock
was used, and how many moves were available. These features have two roles.
They are the inputs of the GBDT baseline, and so define the level that a
network reading every position has to exceed. They can also be given to a
network as additional context features, which tests whether the network
already extracts the same information from the sequence.

\else
\paragraph{Scalar context features.}
Information that does not belong to a square is given as scalar context
features: the four castling rights and the scalar inputs specific to the
task. This task uses \V{task.context.count} (Table~\ref{tab:reads}). Every
model receives them, including the gradient-boosted trees, which never see
the board. For a network they are standardised with the mean and standard
deviation of the training set.

\begin{table}[htbp]\centering\small
\block{reads}
\caption{Scalar context features: the inputs given to every model in addition
to the board.}
\label{tab:reads}
\end{table}

\paragraph{Hand-crafted features.}
Gradient-boosted trees cannot take 64 tokens as input, so we also compute
\V{task.hand.count} hand-crafted features of the position
(Table~\ref{tab:hand} in the appendix): material, mobility, available checks
and captures, hanging pieces, passed pawns and attacks near each king, and,
where the task gives a line of moves, features of that line. These features
have two roles. They are the inputs of the GBDT baseline, and so define the
level that a network reading the board has to exceed. They can also be given
to a network as additional context features, which tests whether the network
already extracts the same information from the board.

\fi

\iflook
\begin{figure}[htbp]\centering
\block{encoding}
\ifsequence
\caption{One position of one example, shown twice: the position before the
player's most costly move of the game. Left: as played, with arrows for the
opponent's preceding move (grey) and the player's move (blue). Right: as
encoded, oriented so that the mover plays up the board, with the squares set
in an auxiliary plane shaded grey for the opponent's move and blue for the
player's move. Every position of the sequence is encoded in this way.}
\else
\caption{One example from the dataset, shown twice. Left: as played, with
arrows for the opponent's preceding move (grey) and the move the example
concerns (blue). Right: as encoded, oriented so that the mover plays up the
board, with the squares set in an auxiliary plane shaded grey for the
opponent's move and, where the task marks a line of moves, blue for the
mover's moves and orange for the replies.}
\fi
\label{fig:encoding}
\end{figure}
\fi


\paragraph{Withholding the solution.}
The dataset records the solution of every puzzle, and the companion task
encodes it. In this task it is withheld from every model. Of the auxiliary
planes, a model receives the en passant target square and the origin and
destination squares of the opponent's preceding move, which the solver also
sees played. The planes that mark the moves of the solution are not used. The
scalar context features are the four castling rights. The number of moves the
solver has to find, which is a context feature in the companion task, is not
given. Of the hand-crafted features, a model receives those that are computed
from the position: material, mobility, hanging pieces, passed pawns, attacks
near each king, and whether the opponent's preceding move was a capture. The
two groups that describe the first move of the solution and the solution line
as a whole are excluded. These restrictions are part of the task definition
and not of a model's configuration, so no model of this task can be
configured to receive the solution.
\iflook
In Figure~\ref{fig:encoding} the first move of the solution is drawn on the
left board for the reader. The encoded position on the right marks only the
opponent's move.
\fi

\subsection{Models and training}
\label{sec:models}

% Models and training: the same sequence of models for every task, so each
% model paper inputs this file under its own heading.
% The prose follows docs/writing.md.
Every task is given the same sequence of models, so that each step of the
sequence measures the contribution of one component. A model is fitted on the
training set, selected on the validation set and scored once on the test set.
Table~\ref{tab:runs} lists the models of this task.

\paragraph{Constant baseline.}
\run{constant} predicts the same value for every example: the mean of the
training targets (for a binary target, the base rate), or their median when
the primary metric is an absolute error. A useful model has to improve on it.

\paragraph{Gradient-boosted decision trees.}
\ifcontextgbdt
\run{gbm-context} is a gradient-boosted decision tree model (GBDT), trained
with LightGBM~\cite{lightgbm} on the scalar context features alone. It never
sees the board, and so measures how much of the target the context already
determines. \run{gbm-hand} adds the hand-crafted features.
\else
\run{gbm-hand} is a gradient-boosted decision tree model (GBDT), trained with
LightGBM~\cite{lightgbm} on the scalar context features and the hand-crafted
features.
\fi
We refer to \run{gbm-hand} as \emph{the GBDT baseline}: every other model is
compared with it on the same test examples. GBDT is trained for at most
\V{method.gbm.rounds} boosting rounds with up to \V{method.gbm.leaves} leaves
per tree and a learning rate of \V{method.gbm.rate}. Each tree is fitted on a
random \V{method.gbm.rows.percent}\% of the examples and
\V{method.gbm.features.percent}\% of the features, and boosting stops when the
validation score has not improved for \V{method.gbm.patience} rounds.

\paragraph{Neural networks.}
A network consists of an encoder, which maps the square tokens, the auxiliary
planes and the context features to a vector, and a prediction head, a
two-layer perceptron that maps this vector to the output. We use two encoder
architectures, so that the results show how much the architecture matters for
the task.

The \emph{transformer} encoder~\cite{attention} treats each square as a token.
The embedding of a token is the sum of a learned embedding of its piece, a
learned embedding of its square and a linear projection of the square's
auxiliary planes. The context features are projected by a small perceptron
and enter as one additional token. Every token attends to every other token
in every layer, so distant squares interact from the first layer on. To give
attention the geometry of the board, each attention head adds a learned bias
that depends on the file and rank offset between two
squares~\cite{relative}: files, ranks, diagonals and knight moves all
correspond to fixed offsets. The layers use pre-layer normalisation, and the
encoder outputs the context token concatenated with the mean of the 64 square
tokens.

The \emph{residual convolutional network} (ResNet) is a stack of residual
blocks of $3\times3$ convolutions over the $8\times8$ board, the architecture
that AlphaZero introduced for chess~\cite{alphazero}. Its input is one binary
plane per token value together with the auxiliary planes. A linear projection
of the context features is added at every square after the first convolution,
and the encoder outputs the mean and the maximum over the squares of the
final feature map. Unlike attention, a convolution relates only neighbouring
squares, so the receptive field grows with depth.

\ifsequence
\paragraph{Sequence model.}
An example of this task is a sequence of positions, so a network has two
stages below its prediction head. One of the two encoders above is applied to
every position of the sequence, with the position's per-move features in
place of the context features, and produces one vector per position. A
two-layer transformer then operates on these vectors in the order in which
the moves were played: each vector is a token, a learned positional embedding
is added to it, and the game-level context features enter as one additional
token. The prediction head receives that token concatenated with the mean of
the position tokens. Sequences in a batch are padded to the length of the
longest, and padded elements are masked: they are not encoded and no token
attends to them, so a prediction does not depend on the other examples in
the batch.

Model sizes are hyperparameters. \run{conv-small} is small enough to train
on a CPU. An example consists of tens of positions, so
\run{transformer-small} and \run{transformer-medium} are trained on a GPU.
\run{transformer-medium-hand} is \run{transformer-medium} with the
hand-crafted features added to its context features.
Table~\ref{tab:settings} in the appendix lists the hyperparameters.
\else
Model sizes are hyperparameters. \run{conv-small} and \run{transformer-small}
are small enough to train on a CPU. \run{transformer-medium} is trained on a
GPU, and \run{transformer-medium-hand} is the same network with the
hand-crafted features added to its context features.
Table~\ref{tab:settings} in the appendix lists the hyperparameters.
\fi

\paragraph{Training.}
All networks are trained in the same way: with AdamW~\cite{adamw}, a learning
rate that increases linearly over the first \V{networks.warmup.percent}\% of
the optimisation steps and then follows a cosine schedule to zero, gradient
clipping, and mixed precision on a GPU. The loss matches the primary metric:
mean squared error for a target scored by RMSE, a Huber loss for a target
scored by MAE, and binary cross-entropy for a probability. A regression
target is standardised with the mean and standard deviation of the training
set. After every epoch the network is evaluated on the validation set. We
keep the weights of the best epoch, and stop training when
\V{networks.patience} consecutive epochs have not improved on it.
\ifmore
A model trained on a larger training set is given proportionally fewer
epochs, and a correspondingly shorter patience
(Section~\ref{sec:ablations}).
\fi

\begin{table}[htbp]\centering\footnotesize
\block{runs}
\caption{The models, best first. \emph{Selected} is the number of trees or
the training epoch chosen on the validation set. \emph{Time} is the total
wall-clock time for training and evaluation.}
\label{tab:runs}
\end{table}


\section{Evaluation Protocol}
\label{sec:protocol}

\paragraph{Metrics.}
The primary metric, fixed before the first full training run, is the RMSE on
the test set in rating points; lower is better. We report three secondary
metrics: the mean absolute error (MAE), the coefficient of determination
($R^2$), and Spearman's rank correlation, which measures only whether the
puzzles are ordered correctly. We also report results separately for
\V{bands.count} bands of the true rating.

% Evaluation protocol: the same for every task. Each model paper names its
% own metrics under its section heading and inputs the rest from here.
% The prose follows docs/writing.md.
\paragraph{Confidence intervals.}
Examples from the same game are not independent, so we resample whole games.
The test games are drawn with replacement \V{method.resamples} times, the
metric is computed each time on the examples of the games drawn, and the
2.5th and 97.5th percentiles of these values form the 95\% confidence
interval (CI) that we report with every score.

\paragraph{Paired comparisons.}
To compare two models we compute both scores on the same resampled games and
take their difference. We call this difference the \emph{gain} of the first
model over the second, and sign it so that a positive gain means that the
first model is better. A gain is \emph{statistically significant} when its
95\% CI excludes zero, and is then set in bold in the tables. Unless stated
otherwise, the second model is the GBDT baseline. Two models tend to do well
and badly on the same games, so a paired gain can be significant even when
the two models' separate intervals overlap. These intervals reflect the
sampling of test games only. Each network was trained with a single random
seed, and the variation between training runs is not included.

\paragraph{Model selection.}
\ifchance
The number of trees, the training epoch whose weights are kept and the
calibration parameters are chosen on the validation set.
\else
The number of trees and the training epoch whose weights are kept are chosen
on the validation set.
\fi
The test set is used once per model, for the scores reported here.

\ifchance
\paragraph{Calibration.}
Before a model is scored, its predicted probabilities are rescaled by Platt
scaling~\cite{platt}: a logistic function of the log-odds whose slope and
intercept are fitted on the validation set. The transformation is monotone,
so it leaves a model's ranking of the examples unchanged and affects only its
log loss and its calibration. What remains of a model's miscalibration after
this step cannot be removed by a linear correction of its log-odds.
\fi

\ifpairs
\paragraph{Acceptance criterion.}
Section~\ref{sec:ablations} evaluates modifications of the models. A
modification is \emph{accepted} only if the modified model is better than the
unmodified one on the validation set and its paired gain on the test set is
statistically significant. A modification that is not accepted is still
reported, as a measurement of its effect.
\fi

\paragraph{Tables and figures.}
A model has the same colour and marker in every figure, and the GBDT baseline
is always drawn in orange. The values plotted in a figure are also given in a
table, in the text or in the appendix. All numbers, tables and figures are
generated from the committed result files by the command that builds the
paper.


\section{Results}
\label{sec:results}

\subsection{Main results}

Table~\ref{tab:leaderboard} and Figure~\ref{fig:scores} give the main
results. The constant baseline predicts \V{floor.constant} for every puzzle
and has an RMSE of \V{floor.score}, the standard deviation of the ratings.
The GBDT baseline, with the castling rights and the \V{task.hand.count}
hand-crafted features of the position, reaches \V{reference.score}, which is
only \V{reference.lift.percent}\% below the constant baseline. Aggregate
properties of the position, such as the material and the numbers of legal
moves, checks and captures, are therefore weak predictors of a puzzle's
difficulty.

\begin{table}[htbp]\centering\small
\block{leaderboard}
\caption{Main results on the \V{data.rows.test} test puzzles, best first.
The gain is the paired difference from \V{task.reference} on the same
resampled games. It is positive when the model is better, and bold when
statistically significant.}
\label{tab:leaderboard}
\end{table}

\begin{figure}[htbp]\centering
\block{scores}
\caption{Test RMSE of every model with its 95\% CI. The vertical line marks
the GBDT baseline.}
\label{fig:scores}
\end{figure}

Compared with the GBDT baseline, \V{networks.ahead.of} are significantly
better (Figure~\ref{fig:gains}). Among the networks of the fixed sequence,
which are trained from random initialisation, the small ResNet is the most
accurate. \run{conv-small}, with \V{run.conv-small.parameters.millions}
million parameters, has a gain of \V{run.conv-small.gain} RMSE points over
the GBDT baseline. The medium transformer, with
\V{run.transformer-medium.parameters.millions} million parameters, gains
\V{run.transformer-medium.gain}, the same transformer with hand-crafted
features \V{run.transformer-medium-hand.gain}, and the small transformer
\V{run.transformer-small.gain}. The intervals of the gains of
\run{conv-small} and \run{transformer-medium-hand} do not overlap
(Table~\ref{tab:leaderboard}).
\ifshown
In the companion task the order is reversed: there both medium transformers
are more accurate than the small ResNet (Table~\ref{tab:shown}).
\fi
The transformers were still improving when their training ended
(Figure~\ref{fig:curves}), so this order may reflect the training schedule
and not only the architectures.

\begin{figure}[htbp]\centering
\block{gains}
\caption{Paired gain in RMSE of each network over the GBDT baseline, with
its 95\% CI. A model whose interval lies entirely to the right of the line is
significantly better than the baseline.}
\label{fig:gains}
\end{figure}

\ifchangestart
The two most accurate models are the medium transformers whose encoder is
pretrained on engine evaluations, with gains of
\V{run.transformer-medium-pretrained.gain} and
\V{run.transformer-medium-hand-pretrained.gain} RMSE points over the GBDT
baseline (Section~\ref{sec:ablations}).
\fi
The best model's gain is \V{best.gain.percent}\% of the GBDT baseline's
error, and \V{best.gain.of-lift}\% of the improvement of the GBDT baseline
over the constant baseline. In this task the best network thus improves on
GBDT by more than GBDT improves on a constant prediction.
\ifshown
The companion task shows the opposite: there GBDT on hand-crafted features of
the solution accounts for most of the improvement, and the networks add a
smaller gain.
\fi

\subsection{Ranking accuracy and absolute error}

The best model orders puzzles moderately well. Its Spearman correlation with
the true rating is \V{best.spearman} and its $R^2$ is \V{best.r2}
(Table~\ref{tab:metrics}), compared with \V{reference.spearman} and
\V{reference.r2} for the GBDT baseline. Its MAE is \V{best.mae} rating
points. Chess players are conventionally grouped into rating classes 200
points wide, so the typical error is larger than one class. A prediction
therefore indicates whether a puzzle is likely to be easy or hard, and does
not place it within a rating class.

\begin{table}[htbp]\centering\small
\block{metrics}
\caption{All metrics for every model on the test set. RMSE and MAE are in
rating points (lower is better); for $R^2$ and Spearman's rank correlation
higher is better. The constant baseline has no rank correlation.}
\label{tab:metrics}
\end{table}

\subsection{Results by rating band}

Table~\ref{tab:bands} gives the RMSE by band of the true rating. Because the
bands are defined by the target, they favour predictions near the mean: the
constant baseline has its smallest error in the middle band and its largest
errors in the extreme bands. Errors should therefore be compared between
models within a band, and not between bands.

\begin{table}[htbp]\centering\small
\block{bands-table}
\caption{Test RMSE by band of the true puzzle rating. The bands are defined
by the target, so values are comparable within a column and not between
columns.}
\label{tab:bands}
\end{table}

The best model has a lower RMSE than the GBDT baseline in
\V{bands.best-ahead.of}. The RMSE of the baseline minus that of the best
model is \V{bands.lead.widest} points for puzzles rated
\V{bands.lead.widest.name} and \V{bands.lead.narrowest} points for puzzles
rated \V{bands.lead.narrowest.name}, where a negative value means that the
baseline has the lower error (Figure~\ref{fig:bands}). The baseline's lower
error in the middle bands is not evidence that it rates these puzzles
better. In the band \V{band.3.name} the constant baseline has an RMSE of
\V{band.3.floor}, lower than that of the GBDT baseline (\V{band.3.reference})
and of the best model (\V{band.3.best}). In this task no model predicts the
rating precisely, and a model whose predictions stay close to the mean rating
is accurate in the middle bands and inaccurate elsewhere. The best model
spreads its predictions more widely. It pays for this in the middle bands and
gains much more in the lowest and the two highest bands.

\begin{figure}[htbp]\centering
\block{bands}
\caption{Difference in RMSE between the GBDT baseline and each network, by
rating band (the differences of the values in Table~\ref{tab:bands}, without
intervals). Above the dashed line the network is better.}
\label{fig:bands}
\end{figure}

\subsection{Training curves}

Figure~\ref{fig:curves} shows the validation RMSE of each network of the
fixed sequence after every epoch. Of these, \V{ladder.at-last-epoch.of}
reached their best validation score in the final epoch, so their training
ended while they were still improving. The scores of the fixed sequence are
therefore lower bounds on what these architectures can achieve. A longer
schedule has not been tested in this task.

\begin{figure}[htbp]\centering
\block{curves}
\caption{Validation RMSE of each network of the fixed sequence after every
epoch. The dashed line is the GBDT baseline's validation RMSE.}
\label{fig:curves}
\end{figure}

\subsection{Feature importance of the GBDT baseline}

The GBDT baseline consists of \V{reference.trees} trees over
\V{reference.features} features. Figure~\ref{fig:leans} shows each feature's
share of the total split gain. The three most important features are
\V{reference.leans.1.feature} (\V{reference.leans.1.share}\% of the gain),
\V{reference.leans.2.feature} (\V{reference.leans.2.share}\%) and
\V{reference.leans.3.feature} (\V{reference.leans.3.share}\%). All of them
describe the position. In the companion task the three most important
features all describe the solution. Split gain shows which features the model
used. It does not identify what makes a puzzle difficult,
because correlated features divide the credit between them arbitrarily.

\begin{figure}[htbp]\centering
\block{leans}
\caption{The ten most important features of the GBDT baseline, by share of
the total split gain.}
\label{fig:leans}
\end{figure}

\ifshown
\section{Comparison with Models Given the Solution}
\label{sec:shown}

Every model in this paper has a counterpart in the companion study. The
counterpart has the same name, is trained from an identical configuration
file on the same training puzzles, and is evaluated on the same test
puzzles. The two models differ only in their inputs. The counterpart also
receives the solution: the auxiliary planes that mark its moves, the number
of moves the solver has to find and, if the model uses hand-crafted features,
the features that describe the first move and the solution line.
Table~\ref{tab:shown} compares the \V{shown.runs} models that both studies
have.

We report two quantities for each model. The \emph{difference} is the RMSE
without the solution minus the RMSE with it, that is, the number of rating
points by which the solution improves the model.
\ifshownpaired
It is computed on the same resampled test games for both models and has a
95\% CI.
\fi
The share \emph{retained} is the model's improvement over the constant
baseline without the solution, as a percentage of its counterpart's
improvement over the constant baseline with the solution. The constant
baseline makes the same prediction in both studies, so both improvements are
measured from the same value.

\begin{table}[htbp]\centering\small
\block{shown}
\caption{Test RMSE of each model without the solution (\emph{Hidden}, this
paper) and of the model of the same name with the solution (\emph{Shown},
the companion study, \V{shown.model}), on the same test puzzles.
\emph{Difference} is the hidden RMSE minus the shown RMSE%
\ifshownpaired
, paired on the same resampled test games and bold when statistically
significant%
\fi
. \emph{Retained} is the improvement of the hidden model over the constant
baseline as a percentage of the improvement of the shown model.}
\label{tab:shown}
\end{table}

\paragraph{GBDT.}
Without the solution the GBDT baseline has an RMSE of
\V{shown.reference.here}, which is \V{shown.reference.here.lift.percent}\%
below the constant baseline. With the solution it has an RMSE of
\V{shown.reference.there}, which is \V{shown.reference.there.lift.percent}\%
below. The difference is \V{shown.reference.worth} rating points, and the
model retains \V{shown.reference.kept.percent}\% of the improvement. GBDT
depends on the solution more than any network does. This agrees with the
feature-group ablation of the companion paper, in which removing the features
of the solution line costs the GBDT baseline more than removing any other
group.

\paragraph{Networks.}
Among the networks the difference is smallest for
\V{shown.networks.least.name} (\V{shown.networks.least.worth} rating points)
and largest for \V{shown.networks.most.name}
(\V{shown.networks.most.worth}). Trained from random initialisation, the
medium transformer retains \V{shown.transformer-medium.kept.percent}\% of
the improvement without hand-crafted features and
\V{shown.transformer-medium-hand.kept.percent}\% with them, and the small
ResNet retains \V{shown.conv-small.kept.percent}\%. The two pretrained
transformers retain \V{shown.transformer-medium-pretrained.kept.percent}\%
and \V{shown.transformer-medium-hand-pretrained.kept.percent}\%.

\paragraph{Best models.}
The best model of this paper, \V{shown.best.here.name}, has an RMSE of
\V{shown.best.here}, and the best model of the companion study,
\V{shown.best.there.name}, has an RMSE of \V{shown.best.there}. Without the
solution, the best model thus achieves \V{shown.best.kept.percent}\% of the
improvement over the constant baseline that is achieved with the solution.
\ifshownpaired
On the same resampled test games the model with the solution is
\V{shown.best.verdict} (gain \V{shown.best.gain} RMSE points; 95\% CI
\V{shown.best.gain.low} to \V{shown.best.gain.high}).
\fi

\paragraph{The solution length alone.}
The companion study includes a model that has no counterpart here:
\V{shown.there.gbm-context.name}, a GBDT that never sees the board and
receives only the castling rights and the number of moves the solver has to
find. Its RMSE is \V{shown.there.gbm-context.score}, which is
\V{shown.there.gbm-context.lift.percent}\% below the constant baseline. It is
\V{shown.there.gbm-context.than-reference} the GBDT baseline of this paper by
\V{shown.there.gbm-context.over-reference} rating points and
\V{shown.there.gbm-context.than-best} the best model of this paper by
\V{shown.there.gbm-context.over-best}. These two values are differences
between test scores and have no confidence interval. For GBDT, this single
property of the solution is more informative than all \V{task.hand.count}
hand-crafted features of the position. A network that reads the board
without the solution is nevertheless more accurate than GBDT with the
solution length.

\paragraph{Interpretation.}
The share retained measures what these models extract from a position, and
not what a position contains. GBDT retains the least. A network retains more:
when the solution is withheld it still receives the whole board. Pretraining
raises the share from \V{shown.transformer-medium.kept.percent}\% to
\V{shown.transformer-medium-pretrained.kept.percent}\% for the medium
transformer, and from \V{shown.transformer-medium-hand.kept.percent}\% to
\V{shown.transformer-medium-hand-pretrained.kept.percent}\% with hand-crafted
features. A better encoder of positions therefore retains more, and the
values in Table~\ref{tab:shown} are lower bounds on the share of the
predictable difficulty that is visible in the position
(Section~\ref{sec:limits}).

\ifshownpaired
\subsection{Subgroups}

We repeat the comparison of the two best models within subgroups of the test
set. The difference is computed within each subgroup as for the whole test
set, as a paired difference over resampled games.

The solution matters more the longer it is. The difference increases with
every additional move, from \V{shown.slice.moves.1.gain} RMSE points for
one-move puzzles to \V{shown.slice.moves.5+.gain} for solutions of five or
more moves (Table~\ref{tab:shown-moves}, Figure~\ref{fig:shown-moves}). It
also increases with the number of quiet moves in the solution, from
\V{shown.slice.quiet.0.gain} for solutions without a quiet move to
\V{shown.slice.quiet.3+.gain} for solutions with three or more. By theme it
is largest for \texttt{\V{shown.slice.themes.most.group}}
(\V{shown.slice.themes.most.gain}) and smallest for
\texttt{\V{shown.slice.themes.least.group}}
(\V{shown.slice.themes.least.gain}; Figure~\ref{fig:shown-themes}).

\begin{table}[htbp]\centering\small
\block{shown-moves}
\caption{The two best models by solution length: the number of test puzzles,
the RMSE of the model without the solution (\emph{Hidden}) and of the model
with it (\emph{Shown}), and their paired difference with its 95\% CI, bold
when statistically significant.}
\label{tab:shown-moves}
\end{table}

\begin{figure}[htbp]\centering
\block{shown-moves-figure}
\caption{Paired gain in RMSE from being given the solution, for the best
model, by solution length, with its 95\% CI.}
\label{fig:shown-moves}
\end{figure}

By rating band the difference is largest for puzzles rated
\V{shown.slice.band.most.group} (\V{shown.slice.band.most.gain}) and smallest
for puzzles rated \V{shown.slice.band.least.group}
(\V{shown.slice.band.least.gain}; Table~\ref{tab:shown-band}). This pattern
is expected of any two models of unequal accuracy, because the bands are
defined by the target: the less accurate model predicts closer to the mean,
which costs little in the middle band and much in the extreme bands.

\begin{table}[htbp]\centering\small
\block{shown-band}
\caption{The two best models by band of the true rating, as in
Table~\ref{tab:shown-moves}.}
\label{tab:shown-band}
\end{table}

\begin{figure}[htbp]\centering
\block{shown-themes-figure}
\caption{Paired gain in RMSE from being given the solution, for the best
model, for the \V{slice.themes.compared} most common themes. A puzzle has
several themes and is counted under each.}
\label{fig:shown-themes}
\end{figure}

The solution thus matters least for one-move puzzles and back-rank mates,
and most for long solutions, for solutions with quiet moves and for the
hardest puzzles. A plausible explanation is that the difficulty of a short
forcing tactic is largely determined by the position, which the network
reads, whereas the difficulty of a long solution depends on its later moves,
which the position does not show.
\fi
\fi

\ifchanges
\section{Ablations}
\label{sec:ablations}

% Ablations: what a modification is, the table of paired comparisons and the
% two figures of validation curves. The same for every task, so a model paper
% inputs this file at the start of its ablation section and then reports its
% own findings. Each part appears once the task has the runs it needs, and the
% list below needs at least one of its items.
% The prose follows docs/writing.md.
\ifpairs
The main results compare complete models. An ablation isolates one factor: it
compares a model with its \emph{base}, a model that differs from it in a
single respect. Table~\ref{tab:pairs} lists every such comparison for this
task, with the gain of the modified model over its base on the validation set
and its paired gain on the test set, computed on the same resampled test
games as in the main results. The acceptance criterion of
Section~\ref{sec:protocol} is applied to these pairs. The modifications are
the following.
\begin{itemize}
\ifchangehand
\item \textbf{Hand-crafted features} (suffix \code{-hand}): the network
receives the hand-crafted features as additional context features.
\fi
\ifchangeepochs
\item \textbf{Longer schedule} (suffix \code{-long}): twice the training
epochs, with the same warm-up fraction and cosine schedule.
\fi
\ifchangestart
\item \textbf{Pretraining} (suffix \code{-pretrained}): the transformer
encoder is initialised from a transformer of the same size that was trained
to predict the engine's win probability for positions from the Lichess
evaluation database~\cite{lichessdb}. The piece and square embeddings and all
transformer layers are copied. The prediction head is initialised randomly,
as are the projections of the auxiliary planes that the pretraining task does
not have and, unless the two tasks have the same context features, the
projection of the context features. Training then follows the standard
schedule.
\fi
\ifchangeattacks
\item \textbf{Attack-aware attention bias} (suffix \code{-attacks}): the
relative position bias encodes where two squares lie and not what stands
between them. With this modification each attention head learns four
additional biases for a pair of squares, applied when the piece on one square
attacks the other, when it is attacked from the other, and for the same two
relations through exactly one intervening piece, as in a pin or a skewer. The
relations are computed from the tokens without search, and the biases are
initialised to zero.
\fi
\ifchangegeometry
\item \textbf{No relative position bias} (suffix \code{-no-geometry}): the
relative position bias is removed. This configuration measures the
contribution of the bias and is not a candidate model.
\fi
\ifchangeconv
\item \textbf{ResNet encoder} (prefix \code{conv-medium}): a ResNet with
squeeze-and-excitation in place of the transformer, trained at the higher
peak learning rate used for the convolutional models.
\fi
\ifchangehybrid
\item \textbf{Hybrid encoder} (prefix \code{hybrid-medium}): a few residual
convolutional blocks followed by transformer layers over the resulting
square representations.
\fi
\ifchangelarger
\item \textbf{Larger transformer} (prefix \code{transformer-large}): more layers and
a larger width, with the training schedule unchanged.
\fi
\ifgames
\item \textbf{Training-set size}: the same model trained on a fraction or a
multiple of the training games (Section~\ref{sec:training-size}).
\fi
\iffeatures
\item \textbf{Feature groups}: a model trained without one group of
hand-crafted features at a time (Section~\ref{sec:feature-groups}).
\fi
\end{itemize}

\longblock{pairs}{Paired comparisons. Each row compares a model with a
baseline, in most rows the model it modifies. \emph{Change} states how the
two differ. The gain is that of the model over the baseline and is positive
when the model is better. On
the validation set it is the difference between the two scores. On the test
set it is the paired gain, with its 95\% CI, and is set in bold when
statistically significant.}{tab:pairs}
\fi

\ifvariants
Figure~\ref{fig:variants} shows the validation curves of each model of the
fixed sequence whose training schedule or initialisation was modified,
together with the modified models. A modified model has the colour of its
base and a broken line: long dashes for a longer schedule and dots for a
pretrained encoder.

\begin{figure}[htbp]\centering
\block{variants}
\caption{Validation curves of a model of the fixed sequence (solid line) and
of the models that modify its training schedule or initialisation (broken
lines), after every epoch. One panel per base model, on a common scale. The
dashed horizontal line is the GBDT baseline's validation score.}
\label{fig:variants}
\end{figure}
\fi

\ifshapes
Figure~\ref{fig:shapes} shows the validation curves of the models that
change the architecture and nothing about the training data, each in its own
panel beside its base.

\begin{figure}[htbp]\centering
\block{shapes}
\caption{Validation curves of a model with a modified architecture (broken
line) and of its base (solid line), after every epoch. One panel per modified
model, on a common scale. The dashed horizontal line is the GBDT baseline's
validation score.}
\label{fig:shapes}
\end{figure}
\fi


\ifchangehand
\subsection{Hand-crafted features}

The hand-crafted features of this task describe only the position, and a
network could in principle compute all of them from the square tokens.
Compared with \run{transformer-medium}, \run{transformer-medium-hand} is
\V{pair.transformer-medium-hand.transformer-medium.valid.verdict} than its
base on the validation set and
\V{pair.transformer-medium-hand.transformer-medium.verdict} on the test set
(gain \V{pair.transformer-medium-hand.transformer-medium.gain} RMSE points;
95\% CI \V{pair.transformer-medium-hand.transformer-medium.gain.low} to
\V{pair.transformer-medium-hand.transformer-medium.gain.high}), so the
modification is \V{pair.transformer-medium-hand.transformer-medium.kept}. The
features help, and their contribution is the smallest in
Table~\ref{tab:pairs}.
\ifshown
In the companion task the same two models also differ in the features of the
solution, so the gain reported there measures a different modification.
\fi
\fi

\ifchangestart
\subsection{Pretraining}

Pretraining the encoder on engine evaluations is the most effective
modification in this task. Without hand-crafted features,
\run{transformer-medium-pretrained} is
\V{pair.transformer-medium-pretrained.transformer-medium.valid.verdict} than
its base on the validation set and
\V{pair.transformer-medium-pretrained.transformer-medium.verdict} on the
test set (gain \V{pair.transformer-medium-pretrained.transformer-medium.gain}
RMSE points; 95\% CI
\V{pair.transformer-medium-pretrained.transformer-medium.gain.low} to
\V{pair.transformer-medium-pretrained.transformer-medium.gain.high}), and
the modification is
\V{pair.transformer-medium-pretrained.transformer-medium.kept}. With
hand-crafted features, \run{transformer-medium-hand-pretrained} is
\V{pair.transformer-medium-hand-pretrained.transformer-medium-hand.valid.verdict}
than its base on the validation set and
\V{pair.transformer-medium-hand-pretrained.transformer-medium-hand.verdict}
on the test set (gain
\V{pair.transformer-medium-hand-pretrained.transformer-medium-hand.gain};
95\% CI
\V{pair.transformer-medium-hand-pretrained.transformer-medium-hand.gain.low}
to
\V{pair.transformer-medium-hand-pretrained.transformer-medium-hand.gain.high}),
and the modification is
\V{pair.transformer-medium-hand-pretrained.transformer-medium-hand.kept}.
The best model of this task is \V{best.name}. Both pretrained transformers
are also more accurate than \run{conv-small}, the best network trained from
random initialisation, and the intervals of their gains over the GBDT
baseline lie entirely above its interval (Table~\ref{tab:leaderboard}).

In this task the context features of the model without hand-crafted features
are the four castling rights, the same as in the pretraining task. The
projection of the context features is therefore copied together with the
encoder, and only the prediction head and the weights of the two planes that
mark the opponent's move are initialised randomly. In the model with
hand-crafted features the projection of the context features is initialised
randomly as well.

\ifshown
Pretraining gains more here than in the companion task. There it lowers the
RMSE of the medium transformer from \V{shown.transformer-medium.there} to
\V{shown.transformer-medium-pretrained.there} without hand-crafted features
and from \V{shown.transformer-medium-hand.there} to
\V{shown.transformer-medium-hand-pretrained.there} with them. Here it lowers
the RMSE from \V{shown.transformer-medium.here} to
\V{shown.transformer-medium-pretrained.here} and from
\V{shown.transformer-medium-hand.here} to
\V{shown.transformer-medium-hand-pretrained.here}, more than twice as much in
both cases. The solution is consequently worth less to a pretrained model:
with hand-crafted features the difference between the two tasks is
\V{shown.transformer-medium-hand-pretrained.worth} rating points for the
pretrained model and \V{shown.transformer-medium-hand.worth} for the model
trained from random initialisation (Table~\ref{tab:shown}). A plausible
explanation is that an engine evaluation depends on the tactics available in
a position, so an encoder trained to predict it has already learned features
of the position that a model without the solution otherwise has to learn from
the puzzles alone.
\fi

With a pretrained encoder the hand-crafted features still contribute. The
pretrained model with hand-crafted features has a test RMSE of
\V{run.transformer-medium-hand-pretrained.score} and the pretrained model
without them one of \V{run.transformer-medium-pretrained.score}, a
difference of \V{run.transformer-medium-hand-pretrained.over-twin} points.
The two pretrained models were not compared directly on resampled games, but
the intervals of their gains over the GBDT baseline do not overlap
(Table~\ref{tab:leaderboard}).

The validation curves of the pretrained models lie below those of their
bases after every epoch, and within a few epochs they pass the scores that
their bases reach at the end of training (Figure~\ref{fig:variants}). The
pretrained model without hand-crafted features reached its best validation
epoch at \V{run.transformer-medium-pretrained.best-epoch} of
\V{run.transformer-medium-pretrained.epochs} and the pretrained model with
them at \V{run.transformer-medium-hand-pretrained.best-epoch} of
\V{run.transformer-medium-hand-pretrained.epochs}. Both curves are nearly
flat over the final epochs.
\fi

% Training-set size: the same ablation for every task, so a model paper
% inputs this file inside its ablation section and then reports its own
% findings. It appears once the task has the runs.
% The prose follows docs/writing.md.
\ifgames
\subsection{Training-set size}
\label{sec:training-size}

To estimate how much a model's accuracy depends on the amount of training
data, we retrain it on a half and on a quarter of the training games and
evaluate it on the full test set (Table~\ref{tab:games},
Figure~\ref{fig:games}). The subsets are selected by a hash of the game
identifier and are nested, so the quarter is contained in the half. A network
is given twice the epochs on half the games and four times the epochs on a
quarter, so the number of optimisation steps is constant and only the amount
of data varies.
\ifmore
We also train on more data. A larger sample of the same source is split in
the same way as the standard dataset, so the validation and test sets are
unchanged and no test game appears in the larger training set. A network
trained on a larger set is given proportionally fewer epochs, so that it
again takes the same number of optimisation steps as its base. Unlike a model
trained on a subset, a model trained on more data is a candidate model, and
the acceptance criterion applies to it.
\fi

\begin{table}[htbp]\centering\footnotesize
\block{study-games}
\caption{Training-set size. Each model is trained on a fraction of the
training games\ifmore, or on a multiple of them from the larger sample,\fi{}
and evaluated on the full test set. The gain is paired with the same model
trained on the standard training set. A negative gain means that the model is
worse, and bold marks statistical significance.}
\label{tab:games}
\end{table}

\begin{figure}[htbp]\centering
\block{study-games-figure}
\caption{Test score as a function of the number of training games. Each step
to the right doubles the training set.}
\label{fig:games}
\end{figure}
\fi


% Feature-group ablation: the same ablation for every task, so a model paper
% inputs this file inside its ablation section and then reports its own
% findings. It appears once the task has the runs.
% The prose follows docs/writing.md.
\iffeatures
\subsection{Feature groups}
\label{sec:feature-groups}

\ifsequence
The hand-crafted features fall into three groups: move accuracy, clock use
and the number of available moves.
\else
The hand-crafted features fall into groups: material, mobility, hanging
pieces, passed pawns and king safety, and, where the task gives a line of
moves, the opponent's last move, the first move of the line and the line as a
whole.
\fi
We retrain a model with one group removed at a time
(Table~\ref{tab:features}, Figure~\ref{fig:features}). A significant loss
shows that the model used the group and that no other input provides the same
information. The absence of a loss has two possible explanations: the group
is uninformative, or the same information is available from another group or
from the other inputs. The ablation therefore measures the marginal
contribution of a group given everything else the model receives, which can
be much smaller than the information the group carries on its own.

\begin{table}[htbp]\centering\footnotesize
\block{study-features}
\caption{Feature-group ablation. Each model is trained on all hand-crafted
features except one group. The gain is paired with the model trained on all
features. A negative gain means that the model is worse without the group,
and bold marks statistical significance. Groups are ordered by the size of
the loss.}
\label{tab:features}
\end{table}

\begin{figure}[htbp]\centering
\block{study-features-figure}
\caption{Paired gain of the GBDT baseline without one feature group over the
GBDT baseline with all features, with 95\% CIs. Left of the vertical line,
the model is worse without the group.}
\label{fig:features}
\end{figure}
\fi

\fi

\section{Error Analysis}
\label{sec:error}

% Error analysis: what the analysis is and how it is computed, the same for
% every task. A model paper inputs this file at the start of its error
% analysis, before the subgroups that are specific to its task.
% The prose follows docs/writing.md.
\iflook
The main results identify the best model. They do not show where a model is
accurate and where it is not, which determines how it can be used. We
therefore analyse the test predictions of every model by subgroup. All
results in this section are computed from the \V{look.test-rows} test
examples (analysis of \V{look.date}).

\paragraph{Subgroups.}
We partition the test set by one variable at a time. For each subgroup we
report the number of test examples, the mean observed target, the mean
prediction of a model and the model's score. Two models are compared:
\V{challenger.name}, the best model other than the GBDT baseline, and the
GBDT baseline itself. The gain of the first over the second is computed
within each subgroup as for the whole test set, as a paired difference over
resampled games (\V{look.resamples} resamples), and is omitted for subgroups
with fewer than \V{look.smallest-compared} test examples. A subgroup whose
interval contains zero provides no evidence in either direction. A numeric
variable is divided into ranges, and a range written $a$-$b$ includes $a$ and
excludes $b$.

\paragraph{Two-way tables.}
For a pair of variables we show the mean observed target and the mean
prediction in each cell, on a common colour scale. Cells with fewer than
\V{grid.smallest-cell} test examples are left empty.

\paragraph{Examples.}
Example \ifsequence games\else positions\fi{} are selected by fixed rules and
not by inspection: typical examples chosen by a hash of their key, the
examples with the best model's most extreme predictions in each direction,
and its largest errors in each direction, \V{examples.count} of each. They
illustrate the quantities under discussion and are not evidence of how often
a pattern occurs.
\else
\emph{This version of the paper was built before the error analysis was run.
The analysis partitions the test set by rating, clock and properties of the
position, measures calibration, and shows example positions. It is produced
by \texttt{lab analyse} on the machine that holds the dataset and the test
predictions of every model, and this section appears once its output is
committed.}
\fi


\iflook
\paragraph{Variables that describe the solution.}
Several of the variables by which we partition the test set describe the
solution: its length, the properties of its first move, the number of its
quiet moves and whether it ends in mate. No model in this paper receives
them. A model's mean prediction in such a subgroup therefore shows whether
the position reveals the property. If it did, the mean prediction would
follow the mean true rating from one subgroup to the next.

\subsection{Distribution of the errors}

The median absolute error of the best model is \V{best.close.median} rating
points, and 90\% of its predictions are within \V{best.close.p90} points of
the true rating. \V{best.close.within.100}\% of the predictions are within
100 points, \V{best.close.within.200}\% within 200 points and
\V{best.close.within.300}\% within 300 points (Table~\ref{tab:closeness}).
For the GBDT baseline the three shares are
\V{reference.close.within.100}\%, \V{reference.close.within.200}\% and
\V{reference.close.within.300}\%. With rating classes 200 points wide, the
best model places fewer than half of the puzzles within one class of their
true rating.

\begin{table}[htbp]\centering\small
\block{closeness}
\caption{Distribution of the absolute errors on the test set: the median,
the 90th percentile, and the share of predictions within 100, 200 and 300
rating points of the true rating.}
\label{tab:closeness}
\end{table}

\subsection{Calibration}

In Figure~\ref{fig:reliability} the test puzzles are divided into
\V{best.reliability.groups} bins of equal size by the model's prediction. In
the lowest bin the mean prediction of the best model is
\V{best.reliability.lowest.said} and the mean true rating is
\V{best.reliability.lowest.happened}. In the highest bin they are
\V{best.reliability.highest.said} and \V{best.reliability.highest.happened}.
The largest difference in any bin is \V{best.reliability.widest-gap} points,
so the mean prediction is within that distance of the mean true rating at
every predicted level. At the two ends the predictions are too extreme: the
lowest predictions are too low and the highest too high.

Predictions for individual puzzles are imprecise. For a given prediction, the
central 80\% of the true ratings span about \V{best.reliability.spread}
points (shaded band in Figure~\ref{fig:reliability}), so a prediction leaves
the rating of an individual puzzle largely undetermined.

\begin{figure}[htbp]\centering
\block{reliability}
\caption{Mean true rating against mean predicted rating. The test puzzles
are divided into bins of equal size by the prediction, and each point is one
bin. The shaded band is the central 80\% of the true ratings in each bin of
the first model. Points on the diagonal are unbiased.}
\label{fig:reliability}
\end{figure}

Conditioning on the true rating instead of the prediction shows strong
regression towards the mean (Figure~\ref{fig:band-bias}). For puzzles with a
true rating below 1000, the mean prediction of \V{challenger.name} minus the
mean true rating is \V{slice.band.-1000.challenger.off} points, and for
puzzles rated 2200 and above it is \V{slice.band.2200+.challenger.off}
points: easy puzzles are rated too hard and hard puzzles too easy. The GBDT
baseline is more biased in both bands (\V{slice.band.-1000.reference.off} and
\V{slice.band.2200+.reference.off}). The best model does recognise the
hardest puzzles as harder than average: its mean prediction for them is
\V{slice.band.2200+.challenger.said}, above the mean rating of the dataset
(\V{label.mean}). Their mean true rating, however, is
\V{slice.band.2200+.happened}.

\begin{figure}[htbp]\centering
\block{slice-band-bias}
\caption{Mean prediction minus mean true rating, by band of the true rating,
for the two models compared. Above the line a model rates the puzzles of the
band as harder than they are, and below it as easier.}
\label{fig:band-bias}
\end{figure}

\subsection{Solution length}

The solution length is strongly associated with the rating, and no model
receives it. The mean true rating rises from
\V{slice.moves.1.happened} for one-move puzzles to
\V{slice.moves.5+.happened} for solutions of five or more moves
(Table~\ref{tab:moves}). The mean prediction of \V{challenger.name} rises
with the solution length as well, but over a smaller range, from
\V{slice.moves.1.challenger.said} to \V{slice.moves.5+.challenger.said}
(Figure~\ref{fig:slice-moves}). Its mean prediction minus the mean true
rating is \V{slice.moves.1.challenger.off} points for one-move puzzles and
\V{slice.moves.5+.challenger.off} points for solutions of five or more
moves. The position therefore reveals part of the solution length.

\begin{table}[htbp]\centering\footnotesize
\block{slice-moves}
\caption{Test-set results by solution length, which no model receives. For
each group: the number of test puzzles, the mean true rating, the mean
prediction and RMSE of each of the two models, and the paired gain of the
first model over the second.}
\label{tab:moves}
\end{table}

\begin{figure}[htbp]\centering
\block{slice-moves-said}
\caption{Mean true rating and mean prediction of the two models compared, by
solution length.}
\label{fig:slice-moves}
\end{figure}

\V{challenger.name} is significantly better than the GBDT baseline in
\V{slice.moves.ahead.of}, and the gain is largest for one-move puzzles
(\V{slice.moves.1.gain} RMSE points). For these puzzles the GBDT baseline
predicts a mean rating of \V{slice.moves.1.reference.said}, far above the
true mean, whereas the network predicts \V{slice.moves.1.challenger.said}.
Almost all one-move puzzles are mates in one: there are
\V{slice.moves.1.rows} one-move puzzles in the test set, and
\V{slice.themes.mateIn1.rows} test puzzles have the theme \texttt{mateIn1},
which by definition has a one-move solution.
We hypothesise that the network recognises from the board that a mate in one
is available, which none of the hand-crafted features states. For longer
solutions the network distinguishes the lengths much less well. Its mean
predictions for solutions of three, four and five or more moves are
\V{slice.moves.3.challenger.said}, \V{slice.moves.4.challenger.said} and
\V{slice.moves.5+.challenger.said}, whereas the mean true ratings are
\V{slice.moves.3.happened}, \V{slice.moves.4.happened} and
\V{slice.moves.5+.happened}.

\subsection{Themes}

Lichess tags every puzzle with themes that describe the tactical motif (for
example a fork, a pin or a discovered attack), the material won, the phase of
the game and the length of the solution. No model receives the themes as
input. We form one group for each of the \V{slice.themes.compared} most
common themes, and a puzzle belongs to the group of every theme it has.
\V{challenger.name} is significantly better than the GBDT baseline in
\V{slice.themes.ahead.of} (Figure~\ref{fig:themes}).

\begin{figure}[htbp]\centering
\block{slice-themes-gain}
\caption{Paired gain in RMSE of the best model other than the GBDT baseline
over the GBDT baseline, for the \V{slice.themes.compared} most common themes.
A puzzle has several themes and is counted under each. A theme whose interval
includes zero provides no evidence in either direction.}
\label{fig:themes}
\end{figure}

The gain is largest for the theme
\texttt{\V{slice.themes.gain.most.group}}
(\V{slice.themes.gain.most} RMSE points) and smallest for
\texttt{\V{slice.themes.gain.least.group}} (\V{slice.themes.gain.least}).
The next largest gains are for \texttt{mateIn1}
(\V{slice.themes.mateIn1.gain}) and \texttt{oneMove}
(\V{slice.themes.oneMove.gain}), which are nearly the same set of puzzles.
The network's advantage over GBDT is thus largest for back-rank mates and
mates in one. The network underestimates most the themes that describe a long or quiet
solution: its mean prediction minus the mean true rating is
\V{slice.themes.veryLong.challenger.off} points for \texttt{veryLong} and
\V{slice.themes.quietMove.challenger.off} points for \texttt{quietMove}.
Table~\ref{tab:slice-themes} in the appendix gives the values behind the
figure.

\subsection{Properties of the solution}

The companion paper shows that several properties of the solution are
associated with large differences in the mean rating. Here we ask whether the
position reveals them. For each property we give the mean true rating of the
two groups, which depends only on the puzzles, and the mean prediction of
\V{challenger.name}.
\begin{itemize}
\item \textbf{Quiet moves.} A quiet move is a move that is neither a capture
nor a check. Puzzles whose solution has no quiet move have a mean rating of
\V{slice.quiet.0.happened}, and puzzles with one quiet move a mean rating of
\V{slice.quiet.1.happened}. The mean predictions are
\V{slice.quiet.0.challenger.said} and \V{slice.quiet.1.challenger.said}
(Table~\ref{tab:quiet}).
\item \textbf{Mate.} Puzzles whose solution ends in mate have a mean rating
of \V{slice.mate.yes.happened}, compared with \V{slice.mate.no.happened} for
puzzles whose solution wins material. The mean predictions are
\V{slice.mate.yes.challenger.said} and \V{slice.mate.no.challenger.said}
(Table~\ref{tab:mate}).
\item \textbf{Check on the first move.} Puzzles whose first move gives check
have a mean rating of \V{slice.check.yes.happened}, compared with
\V{slice.check.no.happened} for the others. The mean predictions are
\V{slice.check.yes.challenger.said} and \V{slice.check.no.challenger.said}.
\item \textbf{Sacrifice on the first move.} When the first move leaves the
moved piece open to capture, the mean rating is
\V{slice.sacrifice.yes.happened}, compared with
\V{slice.sacrifice.no.happened} otherwise. The mean predictions are
\V{slice.sacrifice.yes.challenger.said} and
\V{slice.sacrifice.no.challenger.said}.
\end{itemize}
For each of the four properties the mean predictions of the two groups
differ in the same direction as the mean true ratings, and by less. The
position therefore reveals each property in part, and none of them
completely. The network is significantly better than the GBDT baseline in
every group of these four partitions (Tables~\ref{tab:quiet}
and~\ref{tab:mate}, and Tables~\ref{tab:slice-check}
and~\ref{tab:slice-sacrifice} in the appendix).

\begin{table}[htbp]\centering\footnotesize
\block{slice-quiet}
\caption{Test-set results by the number of quiet moves (neither captures nor
checks) among the solver's moves, which no model receives.}
\label{tab:quiet}
\end{table}

\begin{table}[htbp]\centering\footnotesize
\block{slice-mate}
\caption{Test-set results by whether the solution ends in mate, which no
model receives.}
\label{tab:mate}
\end{table}

\subsection{Reliability of the rating}

As in the companion task, the network's gain over the GBDT baseline depends
on how well a puzzle's rating is measured. By rating deviation, the gain is
\V{slice.deviation.-75.gain} RMSE points for the most reliable ratings (a
deviation below 75) and \V{slice.deviation.90+.gain} for the least reliable
(a deviation of 90 or more), and it increases from each group to the next
(Table~\ref{tab:slice-deviation}). By number of attempts, the gain falls from
each group to the next, from \V{slice.plays.-200.gain} for puzzles attempted
fewer than 200 times to \V{slice.plays.5000+.gain} for puzzles attempted
5{,}000 times or more (Table~\ref{tab:slice-plays}).

These groups also differ in difficulty. The most reliable ratings have a
mean of \V{slice.deviation.-75.happened} and the least reliable a mean of
\V{slice.deviation.90+.happened}, and in this task the gain depends strongly
on the rating band and is negative in the middle bands
(Figure~\ref{fig:bands}). The subgroup tables therefore do not show whether
the reliability of the rating has an effect of its own. An analysis by rating
deviation within each rating band would separate the two effects and has not
been done.

\subsection{Solution length and material}

Figure~\ref{fig:grid} shows the mean rating by solution length and by the
number of minor and major pieces on the board, which serves as a proxy for
the phase of the game. The model receives the pieces and not the solution
length. Its predictions reproduce the direction of the effect of the
solution length and understate its size, as in Table~\ref{tab:moves}.

\begin{figure}[htbp]\centering
\block{grid-moves-pieces}

\medskip
\block{grid-moves-pieces-said}
\caption{Mean rating by solution length and by the number of knights,
bishops, rooks and queens on the board. Top: true ratings. Bottom:
predictions of the best model other than the GBDT baseline. Both panels use
the same colour scale, and cells with too few test puzzles are left empty.}
\label{fig:grid}
\end{figure}

\subsection{Agreement with the GBDT baseline}

Table~\ref{tab:beside} gives, for each model, the correlation of its
predictions with those of the GBDT baseline and the share of test puzzles on
which it is the closer of the two to the true rating. The best model is
closer than the GBDT baseline on \V{best.beside.closer}\% of the test
puzzles, and the correlation between the two sets of predictions is
\V{best.beside.correlation}.

\begin{table}[htbp]\centering\small
\block{beside}
\caption{Agreement with the GBDT baseline on the test set: the correlation
between the two models' predictions, and the share of puzzles on which the
model is closer to the true rating than the baseline.}
\label{tab:beside}
\end{table}

\subsection{Example positions}
\label{sec:positions}

The following examples were selected from the test set by the fixed rules
described above. Each board is shown from the mover's side. The grey arrow is
the opponent's preceding move, which the models receive. The blue arrow is
the first move of the solution, which they do not receive: it is drawn for
the reader. Next to each board we give the solution, the true rating, the
predictions of the best model and of the GBDT baseline, and the puzzle's
identifier on Lichess.

The largest errors show what the position does not reveal. The three largest
underestimates all have the theme \texttt{veryLong}
(Figure~\ref{fig:under}), and each of the three largest overestimates has
the theme \texttt{oneMove} or \texttt{short} (Figure~\ref{fig:over}). In
these puzzles the length of the solution contradicts what the model infers
from the position, which agrees with the results by solution length above.
The puzzles with the highest and the lowest predictions are shown in the
appendix (Figures~\ref{fig:highest} and~\ref{fig:lowest}).

\examples{examples-typical}{Typical puzzles: the first \V{examples.count}
test puzzles in the order of a hash of their key.}{fig:typical}

\examples{examples-over}{The largest overestimates: the test puzzles for
which the best model's prediction exceeds the true rating by the most.}{fig:over}

\examples{examples-under}{The largest underestimates: the test puzzles for
which the best model's prediction falls below the true rating by the
most.}{fig:under}
\fi

\section{Practical Use}
\label{sec:use}

\begin{itemize}
\item \textbf{Where a solution is available, the model that receives it
should be used.} In the coaching application a tactic is rated after an
engine has found it, so its line is known.
\ifshown
With the solution, the best model of the companion study has an RMSE of
\V{shown.best.there} rating points on these test puzzles, compared with
\V{best.score} for the best model without it.
\fi
\item \textbf{Ordering positions before a line is computed.} For a set of
positions that are known to contain a tactic and for which no line has been
computed, the model can order the positions roughly by difficulty: its
Spearman correlation with the true rating is \V{best.spearman}. Its MAE of
\V{best.mae} rating points is too large to state the difficulty of any one of
them.
\item \textbf{Explaining why a tactic was hard.}
\ifshown
The best model achieves \V{shown.best.kept.percent}\% of the improvement
that is achieved with the solution.
\fi
The appearance of the position is thus one part of a tactic's difficulty,
and the properties of the line, such as its length and its quiet moves, are
another. An explanation of a missed tactic should refer to both.
\item \textbf{Not for positions that may contain no tactic.} Every position
in the dataset contains a winning tactic. For an arbitrary position from a
game, a prediction estimates how hard a tactic would be if one existed, and
gives no indication of whether one exists (Section~\ref{sec:limits}).
\item \textbf{Not for display as a number.} A predicted rating is imprecise
for a single position, and it strongly underestimates the hardest tactics.
\end{itemize}

\section{Limitations}
\label{sec:limits}

\begin{itemize}
\item \textbf{The share retained is a lower bound.} It measures what these
models extract from a position, and not all the information that a position
contains. A stronger encoder retains more, as the pretrained models show
(Section~\ref{sec:ablations}). Our results do not show where the share stops
increasing.
\item \textbf{Single training runs.} Each model and its counterpart were
trained once, with one random seed. The confidence intervals reflect the
sampling of test games and not the variation between training runs.
\item \textbf{Few modifications were tested.} The longer schedule, the
architecture ablations and the training-set sizes of the companion study were
not repeated for this task, and most networks of the fixed sequence were
still improving when their training ended.
\item \textbf{Every position contains a winning tactic.} A puzzle is taken
from a position in which the opponent has just made a mistake and a winning
line exists. The models have seen no position without a tactic, which is the
case for most positions of a game.
\item \textbf{Solvers know that a tactic exists.} A puzzle's rating is
determined by solvers who are told that there is something to find. How
hard the same tactic is to notice during a game, without that knowledge, is a
different question, which puzzle ratings do not answer.
\item \textbf{The target is on the Lichess puzzle scale.} As in the
companion study, the rating depends on which solvers attempt the puzzle, and
relating a prediction to a player's game rating requires a mapping that this
work does not provide.
\item \textbf{Only stable ratings are included.} Puzzles with few attempts
are excluded.
\end{itemize}

\section{Reproducibility}
\label{sec:cost}

Training all \V{runs.count} models took \V{cost.seconds} in total:
\V{cost.seconds.networks} for the networks, each trained on a single cloud
GPU, and \V{cost.seconds.others} for the constant baseline and GBDT on a CPU.
The best model has \V{best.parameters} parameters and was trained for
\V{best.epochs.run} epochs in \V{best.seconds} on an \V{best.device}, at
about \V{best.rows-a-second} puzzles per second. Table~\ref{tab:runs} lists
the hardware and training time of every model, and
Table~\ref{tab:settings} the hyperparameters of the networks of the fixed
sequence.

All results are generated from committed files. The dataset is built from
its recipe by \texttt{lab data puzzles}, the models are trained by
\texttt{lab all puzzle-difficulty-unsolved}, each training run writes a
record to the model's \texttt{results/} directory, and
\texttt{lab paper puzzle-difficulty-unsolved} generates every number, table
and figure of this paper from those records, and from the records of the
companion study for Section~\ref{sec:shown}, and compiles it. A record names
the build of the dataset it used (\V{data.fingerprint}) and the commit of
the code that produced it (\V{code.commit}). Model weights and test
predictions are not stored in the repository. The model's card states where
they are kept and gives the hash of the best model's weights.

\section{Conclusion and Future Work}
\label{sec:conclusion}

From the position alone, the rating of a Lichess puzzle can be predicted with
an RMSE of \V{best.score} rating points, compared with \V{floor.score} for a
constant prediction\ifshown{} and \V{shown.best.there} for the same model
given the solution\fi. Aggregate features of the position predict little, a
network that reads the board predicts much more, and pretraining the encoder
on engine evaluations is the largest single improvement. The remaining error
is dominated by properties of the solution that the position reveals only in
part, above all its length. The following experiments are the natural next
steps.
\begin{itemize}
\item \textbf{A stronger pretrained encoder.} The encoder was pretrained
once, at one size, on part of the evaluation database. Pretraining gained
more than any other modification in this task, so longer pretraining, more
positions and a larger encoder are the first experiments to add.
\item \textbf{Attack-aware attention bias.} In the companion task this bias
is the largest gain for a transformer trained from random initialisation.
If part of the gap between the two tasks is due to the network failing to
find the tactic that the position contains, attention that receives the
attack relations should narrow it.
\item \textbf{A longer schedule and more data.} Most networks of the fixed
sequence were still improving when their training ended. In the companion
task both a longer schedule and more training puzzles improved the medium
transformer. Neither has been tested here.
\item \textbf{A second random seed.} Training each model and its counterpart
again would show how much of the difference between them is due to the
solution and how much to the variation between training runs.
\item \textbf{Locating the tactic first.} A separate model of this project
predicts, from the same position, which squares the moves of the solution
leave or reach. Its output could be given to the difficulty model as
additional input planes.
\end{itemize}

\bibliographystyle{unsrt}
\bibliography{../../../papers/references}

\clearpage
\appendix

\section{Hyperparameters and Features}

\begin{table}[htbp]\centering\footnotesize
\block{settings}
\caption{Hyperparameters of the networks of the fixed sequence, from their
configuration files. \emph{LR} is the peak learning rate. Values shared by
all of them are given below the table.}
\label{tab:settings}
\end{table}

\begin{table}[htbp]\centering\footnotesize
\block{hand}
\caption{Hand-crafted features: the inputs of \V{task.reference} in addition
to the scalar context features, and of a network whose configuration includes
them. All of them are computed from the position and the opponent's
preceding move.}
\label{tab:hand}
\end{table}

\iflook
\section{Examples with Extreme Predictions}

\examples{examples-highest}{The test puzzles with the highest predictions of
the best model.}{fig:highest}

\examples{examples-lowest}{The test puzzles with the lowest predictions of
the best model.}{fig:lowest}

\section{Subgroup Tables}

Each table below partitions the test puzzles by one variable. The columns are
the group, the number of test puzzles, the mean true rating, the mean
prediction and RMSE of each of the two models compared, and the paired gain
of the first model over the second with its 95\% CI, in bold when
statistically significant. The variables that describe the solution are not
inputs of any model.

\block{slices}
\fi

\end{document}
